% Four Components of Economics
\subsection{Four Core Principles of Economics}
\textbf{Economics} focuses on scarcity and how it requires individuals, businesses, and governments to make choices.

\textbf{Economic analysis} focuses on individual decision-making.

\textbf{Resource allocation} is the allocation of an economy's scarce resources among alternative uses.

The \textbf{four core principles} act as the main framework for understanding decision-making:
\begin{enumerate}
    \item Cost-Benefit Principle.
    \item Opportunity Cost Principle.
    \item Marginal Principle.
    \item Interdependence Principle.
\end{enumerate}

\subsubsection{Cost-Benefit Principle}

\textbf{Cost-benefit principle} is when a choice is only worth making if the benefits are greater than or equal to the costs. It is important to note that in the case of forgoing the next best alternative, that cost is included.

For example, suppose we have a chocolate bar. The chocolate bar costs \$2. You think the chocolate bar is worth \$3. Since the benefit you get is greater than the money you lose, you should buy it.

In another example, suppose we are going to get dinner at a restaurant with a friend. The dinner costs \$40; however, the value of the meal is \$20. The value of your friend's company is \$40. If you stay at home, the cost of groceries is \$8. In this case, the benefit is \$60, and the cost is \$48. Since the benefit is greater than the cost, you should go to dinner with your friend.

\underline{\textbf{WTP}}

The maximum amount the consumer is willing to give up to get something. Essentially, it is what we think the value (worth) of a product is.

For the same chocolate bar, based on the fact that we think the chocolate bar is worth \$3, that is our \textbf{WTP}.

\underline{\textbf{Economic Surplus (Between Benefits and Total Costs)}}

An \textbf{economic surplus} is the difference between the total benefits and the total costs. It determines how good a choice is (positive surplus is a good choice, negative surplus is not).

For the same chocolate bar, based on the fact that we think the chocolate bar is worth \$3 and costs \$2, the surplus is \$1. This was a good choice.

\underline{\textbf{Framing Effects}}

\textbf{Framing effects} are when the brain can be influenced by rearranging the words in a statement, even though both statements mean the same thing.

Suppose we have a professor who indicates that a company has 6,000 jobs total.

\textit{Example} - ``Save 2,000 jobs.'' This means that 2,000 jobs would remain, and 4,000 jobs would be lost.

\textit{Example} - ``Lose 4,000 jobs.'' This means that 4,000 jobs are gone, and the rest remain. In this case, 2,000 jobs were saved.

\subsubsection{Opportunity Cost}

\textbf{Opportunity cost} is when you choose one thing, you lose the chance to do the next best alternative, as a result of scarcity.

Suppose we have one hour to allocate time, and we have ordered our priorities:

\begin{enumerate}
    \item Homework; \$50 value.
    \item Hanging out with friends; \$11 value.
    \item Watching Netflix; \$10 value.
    \item Taking a nap; \$15 value.
\end{enumerate}

Suppose I choose to do my homework, which is valued at \$50. I choose to do my homework because the benefit of doing it outweighs the cost of forgoing the next best alternative (the highest-valued next best alternative), which is taking a nap.

\underline{\textbf{Scarcity and Trade-Offs}}

\textbf{Scarcity} can be defined as limited resources but unlimited human wants. This will always be the case.

Because scarcity is a factor, every choice we make requires us to give something up. There are limitations in terms of money, time, energy, and attention. Every decision we make involves a trade-off. When we choose one thing, we lose the chance to choose something else.

\textit{Example} - ``Going to school'' vs. ``Working a full-time job.''

We cannot do both at the same time. This means that choosing one means giving up parts of the other.

For example, if you choose to go to school, you give up the money you could have earned from working. You miss out on money you could have made while working full-time. You also pay tuition. This is an opportunity cost, as if you chose to work instead, you would not pay tuition.

Costs like rent and food are not opportunity costs because regardless of the choice we make, we would have to pay them either way. Time spent studying is also not an opportunity cost, as that time would otherwise be spent working full-time.

\textit{Example} - ``Going to Kurdistan'' vs. ``Staying in Canada''

By going to Kurdistan, I give up staying in Canada. By choosing to stay in Canada, I give up seeing family, friends, and culture from Kurdistan.

\textbf{Note:} Whenever we talk about opportunity cost between two things, we should ask: \textit{What do we give up by making one choice over the other?}

\underline{\textbf{Sunk Cost}}

A \textbf{sunk cost} is something that has already been paid for but cannot be recovered.

For instance, suppose you pay \$12 for a movie ticket. If the movie is terrible, you should leave. The \$12 is already gone either way, so it should not influence your choice.

\textbf{Sunk cost fallacy} is the inability to walk away from something that has been heavily invested in, even when giving up is a better idea.\footcite{decisionlab_sunkcost}

\underline{\textbf{Efficiency vs. Equity}}

\textbf{Efficiency} focuses on maximizing total output with given resources.

\textbf{Productive efficiency} is when the economy is getting the maximum output from the resources it has (usually along a PPF).

\textbf{Allocative efficiency} is when the economy produces the combination of goods and services that society values most.

\textbf{Equity} focuses on how income, wealth, and opportunities are distributed across society.

Example - Suppose we have a pie. Efficiency is about making the pie as big as possible, while equity is about sharing the pie in a way that leads to equal outcomes.

\subsubsection{Marginal Principle}

Individuals and firms should make decisions by comparing the \textit{marginal benefits} and \textit{marginal costs} of an action. Within the marginal principle, we should \textbf{continue doing something if the marginal benefit is greater than or equal to the marginal cost.}

Simply put, instead of asking, \textit{``How much should I do of something?''}, you ask:

\textit{``Should I do one more unit of that something?''}

\textbf{Marginal benefit} is the extra benefit you get from doing one more unit of something.

\textbf{Marginal cost} is the extra cost you face from doing one more unit of something.

The \textbf{rational rule (MB = MC)} states that if something is worth doing, you should keep doing it until marginal benefits equal marginal costs.

\subsubsection{Interdependence Principle}

The \textbf{interdependence principle} states that the best choice depends on other choices you make, the choices others make, developments in other markets, and expectations about the future.

Everything is connected, and choices are influenced by the world around you due to limited resources.

\subsubsection{Economic Activities and Markets}

\textbf{Markets} are places where goods, services, and resources are bought and sold through interactions between firms and households.

\underline{\textbf{Product Market}}

A \textbf{product market} is where goods and services are bought and sold.

Firms supply goods and services, while households demand them.

\underline{\textbf{Labour Market}}

A \textbf{labour market} is a market for human skills and effort.

Households supply labour, while firms demand labour.

\underline{\textbf{Capital Market}}

A \textbf{capital market} is a market for financial capital such as savings, loans, and investment funds.

Savers supply funds, while firms borrow funds.

\subsubsection{Positive vs. Normative Economics}

\textbf{Positive economics} focuses on what will happen if a certain action is taken. These statements are testable and based on facts.

\textbf{Normative economics} focuses on what should be done. These statements are based on personal beliefs or value judgments.

\textit{Example} - Increasing labour costs for firms will lead to a smaller labour market. This statement is testable and therefore represents positive economics.

\textit{Example} - The government should implement a federally mandated minimum wage so workers can earn a living income. This statement reflects a value judgment and therefore represents normative economics.